\scanpage{11}
induit une bijection canonique
\[
\widetilde F_M/\{\pm1\}\simeq F_M.\tag*{(22)}
\]
Ceci posé, \inserted{pour} tout $a\in A_M$, il y a exactement deux éléments de $F_M$ qui le contiennent -- donc
\[
(S_M,A_M,F_M)\overset{\mathrm{def}}{=}P_M\tag*{(23)}
\]
est un $2$-polyèdre combinatoire triangulé \inserted{$P_M$}, ayant $S_M$ comme ensemble des sommets, $A_M$ comme ens. d'arêtes, $F_M$ comme ens. des faces.
% slanted
\margin{Attention, il manque encore la condition de \uline{connexité} sur l'ens des arêtes et faces \unsure{incidentes|incidents} à un sommet\ldots}

\uline{Remarque}\quad Supposons $\Lambda$ $\tfrac12$ local, alors pour $x,y\in\widetilde S_M=M^*$, $\exists\,z\in\widetilde S_M$ tel que $(x,z)\in\mathrm{Rep}(M)$, $(y,z)\in\mathrm{Rep}(M)$, ce qui implique que pour $\dot x,\dot y\in S_M$, $\exists\,\dot z\in S_M$ tel que $\{\dot x,\dot z\}\in A_M$, $\{\dot z,\dot y\}\in A_M$. Donc la distance combinatoire de deux \deleted{point} sommets de $P_M$ est au plus égale à $2$, en particulier $P_M$ est un polyèdre triangulé \uline{connexe}.
% archive sheet number: 10
